%% PiFinder_Stellarmate — overview deck, built on the AVVP Beamer template
%% (github.com/apos/beamer_template_AVVP_eV).  Vendored: tex/, fonts/, pics/.
%% Build:  latexmk -lualatex PFSM_overview.tex      (twice, for the agenda bar)
%%   or:   lualatex PFSM_overview  &&  lualatex PFSM_overview
\documentclass[12pt,english,aspectratio=1609]{beamer}
\usepackage{mathptmx}
\usepackage[T1]{fontenc}
\usepackage[utf8]{luainputenc}
\usepackage{amssymb}
\usepackage{graphicx}
\usepackage{setspace}
\PassOptionsToPackage{normalem}{ulem}
\usepackage{ulem}

\makeatletter
\providecommand{\tabularnewline}{\\}
\newcommand{\lyxdot}{.}
\newcommand\makebeamertitle{\frame{\maketitle}}%
\AtBeginDocument{%
  \let\origtableofcontents=\tableofcontents
  \def\tableofcontents{\@ifnextchar[{\origtableofcontents}{\gobbletableofcontents}}
  \def\gobbletableofcontents#1{\origtableofcontents}
}
\input{tex/avvp_beamer_preamble.tex}
\makeatother

\usepackage{babel}

% screenshots live in the repo's docs tree — no copies needed
\graphicspath{%
  {img/}%
  {../pptx/assets/}%
  {../../images/readme/}%
  {../../images/pfinder_lx200/}%
  {../../images/pfinder_simulator/}%
  {../../images/logo/}%
}

\newcommand{\code}[1]{{\ttfamily #1}}
\beamerdefaultoverlayspecification{}   % show everything at once; overview deck, not a talk-reveal

% English deck -> override the template's German auto-labels
\def\AVVPSubAgendaTitle{Where we are}
\def\AVVPMainAgendaTitle{Agenda}

% colour external URLs; leave internal nav links (footer title, agenda bar) at their surrounding colour
\hypersetup{colorlinks=true,urlcolor=AVVPBlue,linkcolor=.,citecolor=.,filecolor=.}

\begin{document}

\AVVPDarkMode
\AVVPFootlineDark
\AVVPBgFillWidth{pics/AVVPBg_dark.png}{0}{0}

\title[PiFinder\_Stellarmate]{PFSM — Plate-solving push-to,\\
on a full imaging rig}
\subtitle{What PFSM adds between the two —\\
for PiFinder and StellarMate developers}
\author[Axel Pospischil]{Axel Pospischil\enskip(github.com/apos)}
\date[github.com/apos]{Community project — not affiliated with PiFinder or StellarMate}

\makebeamertitle

\AVVPAgendaOnlySections   % 3 top-level sections in the footer bar; agenda interstitial per section, not per subsection

\begin{frame}{Agenda}
\tableofcontents{}
\end{frame}

% ============================================================
\section{Introduction}
% ============================================================

\subsection{For newcomers}

\begin{frame}[t]{Two problems, one telescope}
\begin{columns}[T]
\column{6cm}
\textbf{Where am I pointing?}
\begin{itemize}
\item A Dobson or push-to EQ has no encoders.
\item The software has no idea where the tube looks.
\end{itemize}
\column{6cm}
\textbf{Drive the whole rig}
\begin{itemize}
\item Camera, mount, guiding, focus.
\item All need one control layer: INDI, KStars/Ekos, an app.
\end{itemize}
\end{columns}
\end{frame}

\begin{frame}[t]{Each project owns one half}
\begin{columns}[T]
\column{6cm}
\textbf{PiFinder}
\begin{itemize}
\item Global-shutter camera + plate-solver.
\item Live sky position + push-to arrows.
\item Python, open hardware \& software.
\end{itemize}
\column{6cm}
\textbf{StellarMate}
\begin{itemize}
\item INDI — the device-driver standard.
\item KStars / Ekos — capture \& automation.
\item The StellarMate phone / tablet app.
\end{itemize}
\end{columns}
\end{frame}

\begin{frame}[t]{PFSM is the wiring. INDI is the connector.}
\begin{itemize}
\item \textbf{PiFinder} exposes its solved position as an INDI telescope.
\item \textbf{Mount Bridge} couples that to any INDI mount (optional).
\item \textbf{Your rig} — KStars / Ekos, the mount, the app — stays unchanged.
\end{itemize}
\medskip{}
PiFinder is patched in place, never forked. The mount side is generic INDI.
\end{frame}

% ============================================================
\subsection{Project summary \& focus}
% ============================================================

\begin{frame}[t]{What PFSM is}
\begin{itemize}
\item \textbf{A setup script} — installs \& patches PiFinder onto StellarMate OS.
\item \textbf{\code{diffs/*.diff} patches} — applied to a stock checkout via \code{patch(1)}.
\item \textbf{Three INDI drivers} — LX200, Mount Bridge, Simulator.
\item \textbf{A web Control Center} — stdlib-only, port 8765.
\end{itemize}
\end{frame}

\begin{frame}[t]{What PFSM is not}
\begin{itemize}
\item \textbf{Not a fork of PiFinder} — patches re-applied every update, kept minimal.
\item \textbf{Not a fork of StellarMate} — stock INDI, Web Manager, KStars.
\item \textbf{Not a replacement for the INDI Control Panel} — it stays for advanced work.
\item \textbf{Not PiFinder-version-specific} — talks only to the stable \code{/api/*}.
\end{itemize}
\end{frame}

\begin{frame}[t]{Why INDI}
\textbf{Both sides already speak it.}
\begin{itemize}
\item PiFinder exposes its position as an INDI telescope.
\item The Mount Bridge couples that to any INDI mount — no per-mount code.
\end{itemize}
\end{frame}

\begin{frame}[t]{Versions \& platforms}
\textbf{PiFinder 2.6.3 \quad$\cdot$\quad StellarMate OS 2.3.0}
\begin{itemize}
\item Pi 4 — fully tested.
\item Pi 5 — fully tested.
\item UTM x86 — install + Control Center + Mount Bridge vs.\ the Simulator (no real solving).
\end{itemize}
\medskip{}
Pinned to a fixed release tag, not the upstream \emph{release} branch.
\end{frame}

% ============================================================
\section{Implementation}
% ============================================================

\subsection{Control Center}

\begin{frame}[t]{Why a Control Center}
\begin{columns}[T]
\column{0.34\textwidth}
\begin{itemize}
\item Watch an install from a phone at the scope.
\item Switch Fake Mode $\leftrightarrow$ the real service.
\item Check camera / IMU / GPS are really detected.
\item Reboot the Pi without SSH.
\end{itemize}
\column{0.64\textwidth}
\vspace*{-11mm}\hfill\includegraphics[height=0.92\textheight,keepaspectratio]{cc_full_page}
\end{columns}
\end{frame}

\begin{frame}[t]{Stdlib only — by necessity}
\begin{columns}[T]
\column{0.34\textwidth}
\begin{itemize}
\item It bootstraps the very venv that PiFinder needs, so it can't depend on one.
\item \code{http.server} on \code{:8765}, polled every 1--2\,s.
\item HTTP Basic auth against the system account (PAM).
\end{itemize}
\column{0.64\textwidth}
\vspace*{-11mm}\hfill\includegraphics[height=0.92\textheight,keepaspectratio]{cc_full_page}
\end{columns}
\end{frame}

\begin{frame}[t]{Traffic-light badges}
\begin{center}

\includegraphics[width=0.92\linewidth,keepaspectratio]{badge_row}
\end{center}
\begin{itemize}
\item \textbf{White} — not applicable yet.
\item \textbf{Green} — verified functional.
\item \textbf{Yellow} — degraded / estimate.
\item \textbf{Red} — broken or wrong.
\end{itemize}
The rule: verify real state, never ``the process is alive''.
\end{frame}

\begin{frame}[t]{Coupling, without the INDI panel}
\begin{columns}[T]
\column{0.40\textwidth}
\begin{itemize}
\item Quick Actions — one-shots.
\item Coupling presets — one click.
\item Settings — set once, hidden.
\end{itemize}
\column{0.56\textwidth}
\includegraphics[width=\linewidth,height=0.62\textheight,keepaspectratio]{cc_mount_bridge_coupling}
\end{columns}
\end{frame}

% ============================================================
\subsection{PF LX200 \& Mount Bridge}
% ============================================================

\begin{frame}[t]{Three building blocks}
\begin{enumerate}
\item \textbf{PiFinder LX200} — INDI telescope. Reports the solved position; a GoTo becomes a push-to. Required.
\item \textbf{PiFinder Mount Bridge} — optional. Snoops PiFinder + the mount and couples them. Generic INDI only.
\item \textbf{Your mount's driver} — not ours. Only for a motorised mount.
\end{enumerate}
\medskip{}
\begin{center}
\includegraphics[width=0.82\linewidth,keepaspectratio]{mb_diagram}
\end{center}
\end{frame}

\begin{frame}[t]{The Web Manager — a ``Remote'' profile}
\begin{columns}[T]
\column{0.42\textwidth}
\begin{itemize}
\item PFSM's INDI drivers live \textbf{only} in the StellarMate Web Manager (\code{:8624}) — the Equipment Profile is built there.
\item KStars / Ekos then connects in \textbf{Remote Host} mode, never ``Local''.
\item The one real difference from a stock StellarMate setup.
\end{itemize}
\column{0.56\textwidth}
\hfill\includegraphics[width=\linewidth,height=0.78\textheight,keepaspectratio]{webmanager_full}
\end{columns}
\end{frame}

\begin{frame}[t]{One weak PiFinder, one strong box}
Three install roles: \textbf{All-in-one} \quad$\cdot$\quad \textbf{PiFinder host} \quad$\cdot$\quad \textbf{Control host}.
\smallskip{}
\begin{itemize}
\item \textbf{INDI-only install} (\code{--mode=indi\_only}) — just the INDI deps + the two PFSM drivers; no PiFinder clone, no venv.
\item \textbf{Control host} — a strong StellarMate PC runs INDI, Ekos, guiding and capture; the low-power PiFinder Pi over the LAN only supplies its solved position.
\end{itemize}
\vspace{0.4em}
\begin{center}
\includegraphics[width=0.94\linewidth,keepaspectratio]{cc_install}
\end{center}
\end{frame}

\begin{frame}[t]{PiFinder LX200}
\begin{columns}[T]
\column{0.34\textwidth}
\begin{itemize}
\item Talks LX200 to PiFinder's \code{pos\_server.py} (TCP 4030).
\item Reports the solved position each poll; a GoTo is only a push-to.
\item No Sync — PiFinder measures the truth every frame.
\end{itemize}
\column{0.64\textwidth}
\hfill\includegraphics[width=\linewidth,height=0.74\textheight,keepaspectratio]{lx200_cp}
\end{columns}
\end{frame}

\begin{frame}[t]{The Coupling Dial}
\begin{columns}[T]
\column{0.44\textwidth}
\begin{enumerate}
\item \textbf{Off}\\ pure push-to
\item \textbf{Verify / Alert only}\\ warn on drift, hands off
\item \textbf{Auto-correct on drift}\\ Sync past the threshold
\item \textbf{Goto-Forward}\\ real Goto, then settle \& refine
\end{enumerate}
\column{0.54\textwidth}
\includegraphics[width=\linewidth,height=0.66\textheight,keepaspectratio]{cc_map}
\end{columns}
\end{frame}

\begin{frame}[t]{Generic INDI, nothing mount-specific}
Only \code{EQUATORIAL\_EOD\_COORD} + \code{ON\_COORD\_SET} — never a mount command.
\begin{columns}[T]
\column{0.33\textwidth}
\begin{itemize}
\item Multi-Point Alignment (\#191) — automated multi-star run.
\item Reposition detection — adopt or revert a move.
\item Shadow Sync — mirror onto a non-driving device.
\end{itemize}
\column{0.65\textwidth}
\hfill\includegraphics[width=\linewidth,height=0.72\textheight,keepaspectratio]{kstars_menu}
\end{columns}
\end{frame}

% ============================================================
\subsection{PF Simulator}
% ============================================================

\begin{frame}[t]{A mount side and a PiFinder side}
\begin{columns}[T]
\column{0.48\textwidth}
\textbf{Mount side}\\[0.3em]
Stock INDI Telescope Simulator — Sync or GoTo like a real mount.
\vspace{2.6em}

\includegraphics[width=\linewidth,keepaspectratio]{sim_mount_c}
\column{0.48\textwidth}
\textbf{PiFinder side}
\begin{itemize}
\item \textbf{Simulator} — holds a settable RA/Dec.
\item \textbf{Truth Injector} — POSTs to \code{/api/fake\_solve}.
\end{itemize}
\vspace{1.4em}

\includegraphics[width=\linewidth,keepaspectratio]{sim_pf_c}
\end{columns}
\end{frame}

\begin{frame}[t]{Simulator details}
\begin{itemize}
\item \textbf{Sync == Goto} — no motion to model; the RA/Dec just becomes the held value.
\item \textbf{Inject every cycle} — keeps PiFinder's solve timestamp fresh.
\item \textbf{\code{FOLLOW\_MOUNT\_DEVICE}} (\#239) — optionally dead-reckon along a mount's slews.
\end{itemize}
\medskip{}
Next (v3.x): render a GSC star field and run PiFinder's real solver on it — issue \#344.
\end{frame}

% ============================================================
\section{Contribute}
% ============================================================

\subsection{Upstream: PiFinder}

\begin{frame}[t]{PiFinder — patches that aren't PFSM-specific}
\code{docs/upstream\_patch\_inventory.md}
\begin{itemize}
\item Test Mode toggle + status readback — a reliable API trigger.
\item Camera-process resilience — fall back to the debug camera instead of crashing.
\item \code{all\_ips()} — show every network address, not just the default route.
\item LX200 ``no position yet'' — return \code{None}, not a parseable fake coordinate.
\item Near-pole \code{/api/fake\_solve} RA — filed as \code{brickbots/PiFinder\#645}.
\end{itemize}
Ready-to-file PR text: \code{docs/upstream\_pr\_templates.md}.
\end{frame}

\subsection{App: StellarMate}

\begin{frame}[t]{StellarMate — extending the SMOS app for PFSM}
What a thin integration needs:
\begin{itemize}
\item Start / stop the selected INDI profile.
\item Add / remove just the PiFinder drivers.
\item Pick ``the mount'' from the loaded drivers.
\item Three Coupling presets, in PiFinder's words.
\end{itemize}
The coupling logic is already framework-agnostic Python — a thin adapter ports it.\\
Concept: \code{docs/concepts/mount\_bridge\_web\_integration.md}.
\end{frame}

\begin{frame}[t]{CC feature $\rightarrow$ app screen}
\begin{columns}[T]
\column{0.40\textwidth}
\begin{itemize}
\item Status badges $\rightarrow$ a read-only status card.
\item Coupling presets $\rightarrow$ a segmented control.
\item Quick Actions $\rightarrow$ buttons on that card.
\end{itemize}
APIs already exist: \code{/api/mount\_bridge\_*}, or the INDI properties directly.
\column{0.56\textwidth}
\includegraphics[width=\linewidth,height=0.62\textheight,keepaspectratio]{cc_map}
\end{columns}
\end{frame}

% ============================================================
\section*{Summary}
% ============================================================

\begin{frame}[t]{Roadmap \& where to jump in}
\begin{columns}[T]
\column{6cm}
\textbf{v2.x — consolidate}
\begin{itemize}
\item Harden \& test what exists.
\item Send selected patches upstream.
\item A guiding watcher.
\end{itemize}
\column{6cm}
\textbf{v3.x — integrate}
\begin{itemize}
\item Deeper SMOS integration (\#35, \#36).
\item PFSM actions in the PiFinder app (\#43).
\item A ``real'' GSC-solving simulator (\#344).
\end{itemize}
\end{columns}
\vskip1em
\code{github.com/apos/PiFinder\_Stellarmate} \quad$\cdot$\quad GitHub Project 15
\end{frame}

\begin{frame}[t]{Two takeaways}
\begin{enumerate}
\item \textbf{PiFinder devs} — a triaged inventory of patches that belong upstream, PR text already written.
\item \textbf{StellarMate devs} — the SMOS app could drive basic PFSM use with a thin adapter over existing APIs.
\end{enumerate}
\vskip0pt plus.5fill
\begin{itemize}
\item Outlook
\begin{itemize}
\item Real-sky solver simulation (\#344).
\item Selected patches back to \code{brickbots/PiFinder}.
\item App-side coupling controls.
\end{itemize}
\end{itemize}
\end{frame}

\begin{frame}[plain]{}
\begin{center}
\vfill
{\AVVPTitleFont\color{AVVPBlue}\Large Thanks — and clear skies.}\\[1.6em]

\includegraphics[width=0.5\paperwidth,keepaspectratio]{PiFinder-Stellarmate_Wortmarke_Negativ_fuer-dunklen-hg}\\[1.2em]
{\small \textcolor{white}{\textcopyright{} github.com/apos 2026}}\\[0.4em]
{\footnotesize \textcolor{AVVPGrey}{Unofficial community project — not affiliated with StellarMate or PiFinder.}}
\vfill
\end{center}
\end{frame}

\appendix

\section*{Appendix}

\subsection*{Links}
\begin{frame}[t]{Links}
\begin{itemize}
\item \href{https://github.com/apos/PiFinder_Stellarmate}{\code{github.com/apos/PiFinder\_Stellarmate}} — repo, README, docs.
\item \href{https://github.com/apos/PiFinder_Stellarmate/blob/dev/docs/upstream_patch_inventory.md}{\code{docs/upstream\_patch\_inventory.md}} \; / \;
      \href{https://github.com/apos/PiFinder_Stellarmate/blob/dev/docs/upstream_pr_templates.md}{\code{docs/upstream\_pr\_templates.md}}
\item \href{https://github.com/apos/PiFinder_Stellarmate/blob/dev/docs/concepts/mount_bridge_web_integration.md}{\code{docs/concepts/mount\_bridge\_web\_integration.md}}
\item \href{https://github.com/brickbots/PiFinder}{\code{brickbots/PiFinder}} — upstream \quad$\cdot$\quad
      \href{https://github.com/brickbots/PiFinder/issues/645}{issue \#645} (near-pole RA)
\item \href{https://github.com/users/apos/projects/15}{GitHub Project 15} — tracked roadmap.
\end{itemize}
\end{frame}

\end{document}
